\begin{ouverture}[geometrie]{16}{Triangles et quadrilatères}{Partie 4 · Espace et géométrie}
{Figures usuelles · Propriétés · Constructions}
\objectif{Reconnaître et nommer les triangles et quadrilatères usuels.}
\objectif{Utiliser leurs propriétés de côtés et d’angles.}
\objectif{Construire ces figures avec règle, équerre et compas.}
\prerequis{Perpendiculaires, cercle et report de longueurs.}
\end{ouverture}
\pagestyle{otMain}\essentiel\fontsize{11}{15}\selectfont
\phantomsection\label{cours-16-1}
\essSection{1}{Les triangles}
\begin{definition}Un triangle a trois sommets et trois côtés. Le triangle rectangle a un angle droit. Le triangle isocèle a deux côtés de même longueur. Le triangle équilatéral a ses trois côtés de même longueur.\end{definition}
\label{fig-cours-16-1-cours-triangles}\figureGeo{cours-triangles}

\begin{propriete}Dans un triangle isocèle, les deux angles portés par la base sont égaux. Dans un triangle équilatéral, les trois angles sont égaux. Les codes de la figure affirment ces égalités.\end{propriete}
\begin{attention}Un triangle peut cumuler deux natures : rectangle et isocèle à la fois, comme la moitié d’un carré coupé par sa diagonale. Vérifier chaque propriété séparément avec les instruments.\end{attention}
\clearpage
\phantomsection\label{cours-16-2}
\essSection{2}{Construire un triangle}
\begin{savoirfaire}[Construire au compas]Construire un triangle équilatéral de côté 4 cm : tracer [BC] de 4 cm ; ouvrir le compas à 4 cm ; tracer un arc de centre B puis un arc de centre C ; leur croisement donne le troisième sommet A. Le compas reporte les longueurs sans les mesurer.\end{savoirfaire}
\begin{exemple}Pour un triangle isocèle de base 4 cm et de côtés 6 cm : même méthode, compas ouvert à 6 cm pour les deux arcs. Pour un triangle rectangle : tracer d’abord l’angle droit à l’équerre, puis reporter les deux longueurs sur ses côtés.\end{exemple}
\begin{propriete}L’ordre des étapes compte : d’abord le segment connu, ensuite les arcs, enfin les tracés. Garder les arcs de construction visibles : ils racontent la méthode.\end{propriete}
\begin{savoirfaire}[Contrôler sa figure]Après la construction, contrôler : mesurer les côtés à la règle, tester l’angle droit à l’équerre. Une construction se vérifie toujours avec les instruments.\end{savoirfaire}
\clearpage
\phantomsection\label{cours-16-3}
\essSection{3}{Les quadrilatères}
\begin{definition}Un quadrilatère a quatre côtés. Le rectangle a quatre angles droits ; ses côtés opposés sont parallèles et de même longueur. Le losange a quatre côtés de même longueur. Le carré a quatre angles droits et quatre côtés égaux : il est à la fois rectangle et losange.\end{definition}
\label{fig-cours-16-3-cours-quadrilateres}\figureGeo{cours-quadrilateres}

\begin{propriete}Dans un rectangle, un losange ou un carré, les côtés opposés sont parallèles deux à deux. Le trapèze, le pentagone ou l’hexagone seront précisés plus tard : au CM1, les figures de référence sont le triangle, le carré, le rectangle et le losange.\end{propriete}
\begin{attention}Une figure posée sur la pointe reste ce qu’elle est : un carré incliné est toujours un carré. C’est la liste des propriétés, pas l’orientation du dessin, qui décide du nom.\end{attention}
\clearpage
\phantomsection\label{cours-16-4}
\essSection{4}{Construire des quadrilatères}
\begin{savoirfaire}[Le rectangle à l’équerre]Construire un rectangle de 6 cm sur 4 cm : tracer [AB] de 6 cm ; élever en A et en B des perpendiculaires à l’équerre ; reporter 4 cm sur chacune ; joindre les deux points obtenus. Contrôler le dernier angle à l’équerre.\end{savoirfaire}
\begin{exemple}Pour un carré de côté $4{,}5$ cm, même méthode avec quatre longueurs égales. Pour un losange de côté 5 cm : tracer [AB], puis depuis A et depuis un point D choisi, reporter 5 cm au compas pour fermer la figure, en gardant les quatre côtés égaux.\end{exemple}
\begin{savoirfaire}[Suivre un programme]Une figure composée se construit étape par étape, en suivant le programme de construction : repérer les points déjà placés, puis les objets qui s’appuient dessus. Relire chaque phrase après son tracé.\end{savoirfaire}
\begin{attention}À main levée d’abord : un croquis codé prépare la construction et évite d’oublier une contrainte. Les instruments viennent ensuite pour le tracé exact.\end{attention}

\clearpage\rubrique{Découverte}
\begin{blocLivre}{ex-16-1}\begin{exercice}[1]{Reconnaître les triangles}
Classer les quatre triangles dessinés : rectangle, isocèle, équilatéral ou quelconque.
\label{fig-ex-16-1}\figureGeo{tri-1}
\end{exercice}\end{blocLivre}
\begin{blocLivre}{ex-16-2}\begin{exercice}[1]{Coder une figure}
Le triangle MNP vérifie MN $=$ MP $=5$ cm. Recopier la figure et placer les codes qui conviennent. Comment s’appelle ce triangle ?
\label{fig-ex-16-2}\figureGeo{tri-2}
\end{exercice}\end{blocLivre}
\begin{blocLivre}{ex-16-3}\begin{exercice}[1]{Le triangle équilatéral}
Construire un triangle équilatéral ABC de côté 4 cm. Laisser les arcs de construction.
\label{fig-ex-16-3}\figureGeo{tri-3}
\end{exercice}\end{blocLivre}
\begin{blocLivre}{ex-16-4}\begin{exercice}[1]{Reconnaître les quadrilatères}
Nommer chacun des quatre quadrilatères dessinés, en justifiant avec les codes.
\label{fig-ex-16-4}\figureGeo{tri-4}
\end{exercice}\end{blocLivre}
\begin{blocLivre}{ex-16-5}\begin{exercice}[1]{Carré ou losange ?}
La figure montre un quadrilatère aux quatre côtés égaux, posé sur la pointe. Zoé dit : « C’est un carré. » A-t-elle forcément raison ?
\label{fig-ex-16-5}\figureGeo{tri-5}
\end{exercice}\end{blocLivre}

\rubrique{Entraînement}
\begin{blocLivre}{ex-16-6}\begin{exercice}[2]{Le triangle rectangle}
Construire un triangle rectangle en A tel que AB $=5$ cm et AC $=3$ cm. Mesurer BC au millimètre près sur la construction.
\label{fig-ex-16-6}\figureGeo{tri-6}
\end{exercice}\end{blocLivre}
\begin{blocLivre}{ex-16-7}\begin{exercice}[2]{Le triangle isocèle}
Construire un triangle isocèle DEF de base EF $=4$ cm et de côtés DE $=$ DF $=6$ cm.
\label{fig-ex-16-7}\figureGeo{tri-7}
\end{exercice}\end{blocLivre}
\begin{blocLivre}{ex-16-8}\begin{exercice}[2]{Le rectangle aux instruments}
Construire un rectangle ABCD avec AB $=6$ cm et BC $=4$ cm. Contrôler le quatrième angle.
\label{fig-ex-16-8}\figureGeo{tri-8}
\end{exercice}\end{blocLivre}
\begin{blocLivre}{ex-16-9}\begin{exercice}[2]{Le carré penché}
Construire un carré de côté $4{,}5$ cm dont un côté est déjà tracé en biais sur la feuille.
\label{fig-ex-16-9}\figureGeo{tri-9}
\end{exercice}\end{blocLivre}
\begin{blocLivre}{ex-16-10}\begin{exercice}[2]{Le losange au compas}
Construire un losange ABCD de côté 5 cm, en commençant par l’angle en A dessiné à main levée sur le croquis.
\label{fig-ex-16-10}\figureGeo{tri-10}
\end{exercice}\end{blocLivre}

\rubrique{Problèmes}
\begin{blocLivre}{ex-16-11}\begin{exercice}[3]{Une figure composée}
Exécuter le programme : construire un rectangle EFGH avec EF $=6$ cm et FG $=3$ cm ; placer le milieu M du segment [EF] ; tracer les segments [MG] et [MH]. Quelle semble être la nature du triangle MGH ?
\label{fig-ex-16-11}\figureGeo{tri-11}
\end{exercice}\end{blocLivre}
\begin{blocLivre}{ex-16-12}\begin{exercice}[3]{Le programme en désordre}
Remettre dans l’ordre les étapes du programme de construction d’un triangle équilatéral : (a) tracer les arcs ; (b) tracer le segment [BC] de 4 cm ; (c) joindre A à B et à C ; (d) ouvrir le compas à 4 cm.
\label{fig-ex-16-12}\figureGeo{tri-12}
\end{exercice}\end{blocLivre}
\begin{blocLivre}{ex-16-13}\begin{exercice}[3]{La figure qui ment}
Un dessin annonce un rectangle, mais l’équerre révèle qu’un angle n’est pas droit. Que peut-on conclure ? Quel nom reste correct pour cette figure ?
\label{fig-ex-16-13}\figureGeo{tri-13}
\end{exercice}\end{blocLivre}
\begin{blocLivre}{ex-16-14}\begin{exercice}[3]{Le plus grand tour}
Un carré a des côtés de 6 cm ; un triangle équilatéral a des côtés de 6 cm aussi. Lequel a le plus grand périmètre ?
\label{fig-ex-16-14}\figureGeo{tri-14}
\end{exercice}\end{blocLivre}
\begin{blocLivre}{ex-16-15}\begin{exercice}[3]{Les diagonales du rectangle}
Construire un rectangle de 8 cm sur 6 cm, tracer ses deux diagonales, puis les mesurer. Que remarque-t-on ?
\label{fig-ex-16-15}\figureGeo{tri-15}
\end{exercice}\end{blocLivre}

\clearpage
